
\documentclass{article}
\usepackage{colortbl}
\usepackage{makecell}
\usepackage{multirow}
\usepackage{supertabular}

\begin{document}

\newcounter{utterance}

\centering \large Interaction Transcript for game `dond', experiment `coop', episode 1 with mistral{-}large{-}3{-}azure.
\vspace{24pt}

{ \footnotesize  \setcounter{utterance}{1}
\setlength{\tabcolsep}{0pt}
\begin{supertabular}{c@{$\;$}|p{.15\linewidth}@{}p{.15\linewidth}p{.15\linewidth}p{.15\linewidth}p{.15\linewidth}p{.15\linewidth}}
    \# & $\;$A & \multicolumn{4}{c}{Game Master} & $\;\:$B\\
    \hline

    \theutterance \stepcounter{utterance}  
    & & \multicolumn{4}{p{0.6\linewidth}}{
        \cellcolor[rgb]{0.9,0.9,0.9}{
            \makecell[{{p{\linewidth}}}]{
                \texttt{\tiny{[P1$\langle$GM]}}
                \texttt{Mängite koostöö läbirääkimiste mängu, milles peate saavutama teise mängijaga kokkulepet selle kohta, kuidas objektide komplekti teie vahel jagada.} \\
\\ 
\texttt{Reeglid:} \\
\texttt{(a) Teie ja teine mängija saate objektide komplekti. Igaühele antakse ka salajane väärtuse funktsioon, mis näitab, kui palju hindasite iga objekti tüüpi.} \\
\texttt{(b) Vahetate teise mängijaga sõnumeid, et kokku leppida selles, kes saab milliseid objekte. Igaüks võib saata kuni 5 sõnumit või lõpetada mängu varajalt, tehes igal ajal salajase ettepaneku.} \\
\texttt{(c) Teil palutakse esitada salajane ettepanek, mis näitab objekte, mida soovite, vormindatuna ruuduliste sulgudega järgmiselt: "{[}Ettepanek: <arv> <objekti nimi>, <arv> <objekti nimi>, <...>{]}"} \\
\texttt{(d) Kui teie ettepanekud on täiendavad, st ettepaanekute täitmiseks on piisavalt objekte, saab iga mängija skoori vastavalt tema saanud objektide väärtuste summale. Muul juhul saavad mõlemad nullpunkte.} \\
\texttt{(e) Teie eesmärk on maksimeerida oma skoori ja teise mängija skoori summa.} \\
\\ 
\texttt{Alustame.} \\
\\ 
\texttt{Saadaolevate objektide komplekt on:} \\
\\ 
\texttt{2 pastakas, 2 sidrunid, 2 pottid} \\
\\ 
\texttt{Teie salajane väärtuse funktsioon on:} \\
\\ 
\texttt{pastakas: 2, sidrun: 0, pott: 3} \\
\\ 
\texttt{OLULINE: Teie sõnumid, välja arvatud salajane ettepanek, edastatakse otse teisele mängijale, nii et ärge lisage mingit vastust reeglitele ega teksti, mis teie sõnumi kuulutaks. Salajase ettepaneku tegemiseks kasutage määratletud vormingut. Ärge kasutage ruudulisi sulge, kui suhtlesite teise mängijaga, või seda interpreteeritakse teie salajase ettepanekuna.} \\
            }
        }
    }
    & \\ \\

    \theutterance \stepcounter{utterance}  
    & \multicolumn{4}{p{0.6\linewidth}}{
        \cellcolor[rgb]{0.9,0.9,0.9}{
            \makecell[{{p{\linewidth}}}]{
                \texttt{\tiny{[P1$\rangle$GM]}}
                \texttt{Ma alustan läbirääkimisi. Mis on teie eelistused või soovid objektide jagamiseks?} \\
            }
        }
    }
    & & \\ \\

    \theutterance \stepcounter{utterance}  
    & & & \multicolumn{2}{p{0.3\linewidth}}{
        \cellcolor[rgb]{0.9,0.9,0.9}{
            \makecell[{{p{\linewidth}}}]{
                \texttt{\tiny{[GM$|$GM]}}
                \texttt{continue} \\
            }
        }
    }
    & & \\ \\

    \theutterance \stepcounter{utterance}  
    & & & \multicolumn{2}{p{0.3\linewidth}}{
        \cellcolor[rgb]{0.9,0.9,0.9}{
            \makecell[{{p{\linewidth}}}]{
                \texttt{\tiny{[GM$|$GM]}}
                \texttt{Ma alustan läbirääkimisi. Mis on teie eelistused või soovid objektide jagamiseks?} \\
            }
        }
    }
    & & \\ \\

    \theutterance \stepcounter{utterance}  
    & & \multicolumn{4}{p{0.6\linewidth}}{
        \cellcolor[rgb]{0.9,0.9,0.9}{
            \makecell[{{p{\linewidth}}}]{
                \texttt{\tiny{[GM$\rangle$P2]}}
                \texttt{Mängite koostöö läbirääkimiste mängu, milles peate saavutama teise mängijaga kokkulepet selle kohta, kuidas objektide komplekti teie vahel jagada.} \\
\\ 
\texttt{Reeglid:} \\
\texttt{(a) Teie ja teine mängija saate objektide komplekti. Igaühele antakse ka salajane väärtuse funktsioon, mis näitab, kui palju hindasite iga objekti tüüpi.} \\
\texttt{(b) Vahetate teise mängijaga sõnumeid, et kokku leppida selles, kes saab milliseid objekte. Igaüks võib saata kuni 5 sõnumit või lõpetada mängu varajalt, tehes igal ajal salajase ettepaneku.} \\
\texttt{(c) Teil palutakse esitada salajane ettepanek, mis näitab objekte, mida soovite, vormindatuna ruuduliste sulgudega järgmiselt: "{[}Ettepanek: <arv> <objekti nimi>, <arv> <objekti nimi>, <...>{]}"} \\
\texttt{(d) Kui teie ettepanekud on täiendavad, st ettepaanekute täitmiseks on piisavalt objekte, saab iga mängija skoori vastavalt tema saanud objektide väärtuste summale. Muul juhul saavad mõlemad nullpunkte.} \\
\texttt{(e) Teie eesmärk on maksimeerida oma skoori ja teise mängija skoori summa.} \\
\\ 
\texttt{Alustame.} \\
\\ 
\texttt{Saadaolevate objektide komplekt on:} \\
\\ 
\texttt{2 pastakas, 2 sidrunid, 2 pottid} \\
\\ 
\texttt{Teie salajane väärtuse funktsioon on:} \\
\\ 
\texttt{pastakas: 3, sidrun: 1, pott: 1} \\
\\ 
\texttt{OLULINE: Teie sõnumid, välja arvatud salajane ettepanek, edastatakse otse teisele mängijale, nii et ärge lisage mingit vastust reeglitele ega teksti, mis teie sõnumi kuulutaks. Salajase ettepaneku tegemiseks kasutage määratletud vormingut. Ärge kasutage ruudulisi sulge, kui suhtlesite teise mängijaga, või seda interpreteeritakse teie salajase ettepanekuna.} \\
\\ 
\texttt{Ma alustan läbirääkimisi. Mis on teie eelistused või soovid objektide jagamiseks?} \\
            }
        }
    }
    & \\ \\

    \theutterance \stepcounter{utterance}  
    & & & \multicolumn{4}{p{0.6\linewidth}}{
        \cellcolor[rgb]{0.9,0.9,0.9}{
            \makecell[{{p{\linewidth}}}]{
                \texttt{\tiny{[GM$\langle$P2]}}
                \texttt{Ma sooviksin saada mõlemad pastakad, sest need on minu jaoks kõige väärtuslikumad. Kui see pole võimalik, olen valmis ka kompromissile. Millised on teie eelistused?} \\
            }
        }
    }
     \\ \\

    \theutterance \stepcounter{utterance}  
    & & & \multicolumn{2}{p{0.3\linewidth}}{
        \cellcolor[rgb]{0.9,0.9,0.9}{
            \makecell[{{p{\linewidth}}}]{
                \texttt{\tiny{[GM$|$GM]}}
                \texttt{continue} \\
            }
        }
    }
    & & \\ \\

    \theutterance \stepcounter{utterance}  
    & & & \multicolumn{2}{p{0.3\linewidth}}{
        \cellcolor[rgb]{0.9,0.9,0.9}{
            \makecell[{{p{\linewidth}}}]{
                \texttt{\tiny{[GM$|$GM]}}
                \texttt{Ma sooviksin saada mõlemad pastakad, sest need on minu jaoks kõige väärtuslikumad. Kui see pole võimalik, olen valmis ka kompromissile. Millised on teie eelistused?} \\
            }
        }
    }
    & & \\ \\

    \theutterance \stepcounter{utterance}  
    & & \multicolumn{4}{p{0.6\linewidth}}{
        \cellcolor[rgb]{0.9,0.9,0.9}{
            \makecell[{{p{\linewidth}}}]{
                \texttt{\tiny{[P1$\langle$GM]}}
                \texttt{Ma sooviksin saada mõlemad pastakad, sest need on minu jaoks kõige väärtuslikumad. Kui see pole võimalik, olen valmis ka kompromissile. Millised on teie eelistused?} \\
            }
        }
    }
    & \\ \\

    \theutterance \stepcounter{utterance}  
    & \multicolumn{4}{p{0.6\linewidth}}{
        \cellcolor[rgb]{0.9,0.9,0.9}{
            \makecell[{{p{\linewidth}}}]{
                \texttt{\tiny{[P1$\rangle$GM]}}
                \texttt{Ma olen huvitatud saama mõlemad potid, sest need on minu jaoks kõige väärtuslikumad. Kui see ei ole võimalik, võiksin kaaluda ühe poti ja ühe pastaka võtmist. Kas teie jaoks oleks see sobiv kompromiss?} \\
            }
        }
    }
    & & \\ \\

    \theutterance \stepcounter{utterance}  
    & & & \multicolumn{2}{p{0.3\linewidth}}{
        \cellcolor[rgb]{0.9,0.9,0.9}{
            \makecell[{{p{\linewidth}}}]{
                \texttt{\tiny{[GM$|$GM]}}
                \texttt{continue} \\
            }
        }
    }
    & & \\ \\

    \theutterance \stepcounter{utterance}  
    & & & \multicolumn{2}{p{0.3\linewidth}}{
        \cellcolor[rgb]{0.9,0.9,0.9}{
            \makecell[{{p{\linewidth}}}]{
                \texttt{\tiny{[GM$|$GM]}}
                \texttt{Ma olen huvitatud saama mõlemad potid, sest need on minu jaoks kõige väärtuslikumad. Kui see ei ole võimalik, võiksin kaaluda ühe poti ja ühe pastaka võtmist. Kas teie jaoks oleks see sobiv kompromiss?} \\
            }
        }
    }
    & & \\ \\

    \theutterance \stepcounter{utterance}  
    & & \multicolumn{4}{p{0.6\linewidth}}{
        \cellcolor[rgb]{0.9,0.9,0.9}{
            \makecell[{{p{\linewidth}}}]{
                \texttt{\tiny{[GM$\rangle$P2]}}
                \texttt{Ma olen huvitatud saama mõlemad potid, sest need on minu jaoks kõige väärtuslikumad. Kui see ei ole võimalik, võiksin kaaluda ühe poti ja ühe pastaka võtmist. Kas teie jaoks oleks see sobiv kompromiss?} \\
            }
        }
    }
    & \\ \\

    \theutterance \stepcounter{utterance}  
    & & & \multicolumn{4}{p{0.6\linewidth}}{
        \cellcolor[rgb]{0.9,0.9,0.9}{
            \makecell[{{p{\linewidth}}}]{
                \texttt{\tiny{[GM$\langle$P2]}}
                \texttt{Ma saan aru. Kui te võtaksite ühe poti ja ühe pastaka, siis ma võiksin võtta ühe pastaka ja ühe sidruni. Kas see teile sobib?} \\
            }
        }
    }
     \\ \\

    \theutterance \stepcounter{utterance}  
    & & & \multicolumn{2}{p{0.3\linewidth}}{
        \cellcolor[rgb]{0.9,0.9,0.9}{
            \makecell[{{p{\linewidth}}}]{
                \texttt{\tiny{[GM$|$GM]}}
                \texttt{continue} \\
            }
        }
    }
    & & \\ \\

    \theutterance \stepcounter{utterance}  
    & & & \multicolumn{2}{p{0.3\linewidth}}{
        \cellcolor[rgb]{0.9,0.9,0.9}{
            \makecell[{{p{\linewidth}}}]{
                \texttt{\tiny{[GM$|$GM]}}
                \texttt{Ma saan aru. Kui te võtaksite ühe poti ja ühe pastaka, siis ma võiksin võtta ühe pastaka ja ühe sidruni. Kas see teile sobib?} \\
            }
        }
    }
    & & \\ \\

    \theutterance \stepcounter{utterance}  
    & & \multicolumn{4}{p{0.6\linewidth}}{
        \cellcolor[rgb]{0.9,0.9,0.9}{
            \makecell[{{p{\linewidth}}}]{
                \texttt{\tiny{[P1$\langle$GM]}}
                \texttt{Ma saan aru. Kui te võtaksite ühe poti ja ühe pastaka, siis ma võiksin võtta ühe pastaka ja ühe sidruni. Kas see teile sobib?} \\
            }
        }
    }
    & \\ \\

    \theutterance \stepcounter{utterance}  
    & \multicolumn{4}{p{0.6\linewidth}}{
        \cellcolor[rgb]{0.9,0.9,0.9}{
            \makecell[{{p{\linewidth}}}]{
                \texttt{\tiny{[P1$\rangle$GM]}}
                \texttt{See teeb kokku:} \\
\texttt{{-} Mina: 1 pott (väärtus 3) + 1 pastakas (väärtus 2) = **5 punkti**} \\
\texttt{{-} Teie: 1 pastakas (väärtus 2) + 1 sidrun (väärtus 0) = **2 punkti**} \\
\\ 
\texttt{Kokku **7 punkti**, mis on optimaalne summa (2 pastakat=4, 2 sidrunit=0, 2 potti=6 → maksimaalne võimalik summa on 10, kuid meie jaotus annab 7/10).} \\
\\ 
\texttt{Alternatiivina võiksin pakkuda:} \\
\texttt{{-} Mina: **2 potti** (väärtus 6)} \\
\texttt{{-} Teie: **2 pastakat** (väärtus 4)} \\
\texttt{Kokku **10 punkti**, mis on maksimaalne võimalik. Kas see teile sobiks?} \\
            }
        }
    }
    & & \\ \\

    \theutterance \stepcounter{utterance}  
    & & & \multicolumn{2}{p{0.3\linewidth}}{
        \cellcolor[rgb]{0.9,0.9,0.9}{
            \makecell[{{p{\linewidth}}}]{
                \texttt{\tiny{[GM$|$GM]}}
                \texttt{continue} \\
            }
        }
    }
    & & \\ \\

    \theutterance \stepcounter{utterance}  
    & & & \multicolumn{2}{p{0.3\linewidth}}{
        \cellcolor[rgb]{0.9,0.9,0.9}{
            \makecell[{{p{\linewidth}}}]{
                \texttt{\tiny{[GM$|$GM]}}
                \texttt{See teeb kokku:} \\
\texttt{{-} Mina: 1 pott (väärtus 3) + 1 pastakas (väärtus 2) = **5 punkti**} \\
\texttt{{-} Teie: 1 pastakas (väärtus 2) + 1 sidrun (väärtus 0) = **2 punkti**} \\
\\ 
\texttt{Kokku **7 punkti**, mis on optimaalne summa (2 pastakat=4, 2 sidrunit=0, 2 potti=6 → maksimaalne võimalik summa on 10, kuid meie jaotus annab 7/10).} \\
\\ 
\texttt{Alternatiivina võiksin pakkuda:} \\
\texttt{{-} Mina: **2 potti** (väärtus 6)} \\
\texttt{{-} Teie: **2 pastakat** (väärtus 4)} \\
\texttt{Kokku **10 punkti**, mis on maksimaalne võimalik. Kas see teile sobiks?} \\
            }
        }
    }
    & & \\ \\

    \theutterance \stepcounter{utterance}  
    & & \multicolumn{4}{p{0.6\linewidth}}{
        \cellcolor[rgb]{0.9,0.9,0.9}{
            \makecell[{{p{\linewidth}}}]{
                \texttt{\tiny{[GM$\rangle$P2]}}
                \texttt{See teeb kokku:} \\
\texttt{{-} Mina: 1 pott (väärtus 3) + 1 pastakas (väärtus 2) = **5 punkti**} \\
\texttt{{-} Teie: 1 pastakas (väärtus 2) + 1 sidrun (väärtus 0) = **2 punkti**} \\
\\ 
\texttt{Kokku **7 punkti**, mis on optimaalne summa (2 pastakat=4, 2 sidrunit=0, 2 potti=6 → maksimaalne võimalik summa on 10, kuid meie jaotus annab 7/10).} \\
\\ 
\texttt{Alternatiivina võiksin pakkuda:} \\
\texttt{{-} Mina: **2 potti** (väärtus 6)} \\
\texttt{{-} Teie: **2 pastakat** (väärtus 4)} \\
\texttt{Kokku **10 punkti**, mis on maksimaalne võimalik. Kas see teile sobiks?} \\
            }
        }
    }
    & \\ \\

    \theutterance \stepcounter{utterance}  
    & & & \multicolumn{4}{p{0.6\linewidth}}{
        \cellcolor[rgb]{0.9,0.9,0.9}{
            \makecell[{{p{\linewidth}}}]{
                \texttt{\tiny{[GM$\langle$P2]}}
                \texttt{{[}Ettepanek: 2 pastakas{]}} \\
            }
        }
    }
     \\ \\

    \theutterance \stepcounter{utterance}  
    & & & \multicolumn{2}{p{0.3\linewidth}}{
        \cellcolor[rgb]{0.9,0.9,0.9}{
            \makecell[{{p{\linewidth}}}]{
                \texttt{\tiny{[GM$|$GM]}}
                \texttt{proposal} \\
            }
        }
    }
    & & \\ \\

    \theutterance \stepcounter{utterance}  
    & & & \multicolumn{2}{p{0.3\linewidth}}{
        \cellcolor[rgb]{0.9,0.9,0.9}{
            \makecell[{{p{\linewidth}}}]{
                \texttt{\tiny{[GM$|$GM]}}
                [2, 0, 0]
            }
        }
    }
    & & \\ \\

    \theutterance \stepcounter{utterance}  
    & & \multicolumn{4}{p{0.6\linewidth}}{
        \cellcolor[rgb]{0.9,0.9,0.9}{
            \makecell[{{p{\linewidth}}}]{
                \texttt{\tiny{[P1$\langle$GM]}}
                \texttt{Teine mängija on esitanud salajase ettepaneku. Nüüd on teie kord esitada oma ettepanek.} \\
            }
        }
    }
    & \\ \\

    \theutterance \stepcounter{utterance}  
    & \multicolumn{4}{p{0.6\linewidth}}{
        \cellcolor[rgb]{0.9,0.9,0.9}{
            \makecell[{{p{\linewidth}}}]{
                \texttt{\tiny{[P1$\rangle$GM]}}
                \texttt{{[}Ettepanek: 2 pott, 0 pastakas, 0 sidrun{]}} \\
            }
        }
    }
    & & \\ \\

    \theutterance \stepcounter{utterance}  
    & & & \multicolumn{2}{p{0.3\linewidth}}{
        \cellcolor[rgb]{0.9,0.9,0.9}{
            \makecell[{{p{\linewidth}}}]{
                \texttt{\tiny{[GM$|$GM]}}
                \texttt{proposal} \\
            }
        }
    }
    & & \\ \\

    \theutterance \stepcounter{utterance}  
    & & & \multicolumn{2}{p{0.3\linewidth}}{
        \cellcolor[rgb]{0.9,0.9,0.9}{
            \makecell[{{p{\linewidth}}}]{
                \texttt{\tiny{[GM$|$GM]}}
                [0, 0, 2]
            }
        }
    }
    & & \\ \\

    \theutterance \stepcounter{utterance}  
    & & & \multicolumn{2}{p{0.3\linewidth}}{
        \cellcolor[rgb]{0.9,0.9,0.9}{
            \makecell[{{p{\linewidth}}}]{
                \texttt{\tiny{[GM$|$GM]}}
                [[0, 0, 2], [2, 0, 0]]
            }
        }
    }
    & & \\ \\

\end{supertabular}
}

\end{document}
